
\documentclass[11pt]{article}
\usepackage[utf8]{inputenc}
\usepackage[spanish]{babel}
\usepackage{amsmath}
\usepackage{geometry}
\geometry{margin=1in}

\title{Documentación: Amazon Series de Tiempo}
\date{}

\begin{document}
\maketitle

\section{Objetivo}
Analizar el comportamiento temporal de las acciones de Amazon mediante herramientas de series de tiempo.

\section{Temas Cubiertos}
\begin{itemize}
\item Descarga y preparación de datos financieros.
\item Visualización de precios históricos.
\item Rendimientos y volatilidad.
\item Estacionariedad.
\item ACF y PACF.
\item Modelado de series temporales.
\end{itemize}

\section{Importancia}
El análisis de series de tiempo permite identificar patrones, tendencias y riesgos asociados al comportamiento histórico de Amazon.

\section{Conclusión}
Las técnicas estudiadas constituyen la base para construir modelos predictivos y de gestión de riesgo financiero.

\end{document}
